\section{Lecture 18 (February 26, 2025, est.): Products, equivariant objects, and tensor powers}
\label{lec18}
\noindent\textit{Scribe: Guillermo Sanmarco}

Throughout this lecture, $k$ is algebraically closed. Our aim is to
construct the Frobenius functor for a semisimple symmetric tensor
category in characteristic $p$. We first make precise the product of
categories and the category of objects equipped with a group action.
These constructions will also explain the product decomposition of
$\operatorname{Ver}_p$ from the previous lecture.

\subsection{Products of linear categories}

Let $\mathcal C$ and $\mathcal D$ be $k$-linear additive Karoubian
categories. Begin with formal pairs $X\boxtimes Y$ and put
\[
 \Hom(X\boxtimes Y,X'\boxtimes Y')
 =\Hom_{\mathcal C}(X,X')\otimes_k
   \Hom_{\mathcal D}(Y,Y').
\]
Composition is componentwise on elementary tensors. One must then
take the additive hull, allowing finite formal direct sums, and its
idempotent completion. We denote the resulting Karoubian product by
$\mathcal C\boxtimes\mathcal D$ in this lecture. Merely taking pairs
would not in general give an additive category.

If the factors are rigid symmetric monoidal, the formulas
\[
 (X\boxtimes Y)\otimes(X'\boxtimes Y')
 =(X\otimes X')\boxtimes(Y\otimes Y'),\qquad
 (X\boxtimes Y)^*=X^*\boxtimes Y^*
\]
extend to give the product the same structures. There are canonical
functors $X\mapsto X\boxtimes\1$ and $Y\mapsto\1\boxtimes Y$.

\begin{proposition}[Universal property]
For a $k$-linear additive Karoubian symmetric monoidal category
$\mathcal A$, giving an additive $k$-linear symmetric monoidal functor
$\mathcal C\boxtimes\mathcal D\to\mathcal A$ is equivalent to giving
one such functor from each factor to $\mathcal A$.
\end{proposition}
\begin{proof}
Given $F:\mathcal C\to\mathcal A$ and $G:\mathcal D\to\mathcal A$,
define $H(X\boxtimes Y)=F(X)\otimes G(Y)$ and
$H(f\otimes g)=F(f)\otimes G(g)$. Its tensor constraint uses the
symmetry to interchange the middle two factors:
\[
 F(X)\otimes G(Y)\otimes F(X')\otimes G(Y')
 \longrightarrow F(X\otimes X')\otimes G(Y\otimes Y').
\]
The symmetry axioms give the required coherence. Additivity and
splitting idempotents extend $H$ to the completed product, uniquely
up to natural isomorphism. Conversely, restrict a functor $H$ to the
two canonical copies of the factors. Since
$X\boxtimes Y=(X\boxtimes\1)\otimes(\1\boxtimes Y)$, the two
procedures are inverse, including on monoidal natural isomorphisms.
\end{proof}

\begin{remark}
This Karoubian product should not be confused with the general abelian
Deligne tensor product, characterized by right-exact bilinear functors.
The distinction disappears in the case used here, where one factor
is split semisimple and the other is abelian. If $\mathcal C$ is split
semisimple with simple objects $S$, then, as a linear category,
\[
 \mathcal C\boxtimes\mathcal D
 \simeq\bigoplus_{S\in\operatorname{Irr}(\mathcal C)}\mathcal D,
\]
where objects have finite support. Indeed, every object of the
product is a finite sum of objects $S\boxtimes Y_S$, there are no
morphisms between distinct $S$-blocks, and each block has exactly the
morphisms of $\mathcal D$. Splitting idempotents preserves this
description. In particular the product is abelian if $\mathcal D$ is.
When both factors are split semisimple, its simples are precisely
the pairs of simples.
\end{remark}

\begin{proposition}\label{prop:lec18-disjoint-product}
Let $\mathcal A$ and $\mathcal B$ be full semisimple tensor
subcategories of a semisimple symmetric tensor category $\mathcal C$.
If their only common simple object is $\1$, then multiplication gives
a symmetric equivalence
\[
 \mathcal A\boxtimes\mathcal B
 \simeq\langle\mathcal A,\mathcal B\rangle,
\]
where the right-hand side is the full tensor subcategory generated by
the two subcategories.
\end{proposition}
\begin{proof}
For simple $X,X'\in\mathcal A$ and $Y,Y'\in\mathcal B$, rigidity
and symmetry identify
\[
 \Hom(X\otimes Y,X'\otimes Y')
 \cong\Hom(X'^*\otimes X,Y'\otimes Y^*).
\]
An object of $\mathcal A$ and an object of $\mathcal B$ can have only
copies of the unit as common simple summands. The last Hom space is
therefore the tensor product of the corresponding unit multiplicity
spaces. Its dimension is $1$ if $X\cong X'$ and $Y\cong Y'$, and
is $0$ otherwise. In particular, the $X\otimes Y$ are simple and
distinct for distinct pairs. This proves full faithfulness of
multiplication on the product, and its image is closed under sums
and summands. Every tensor word in the two subcategories can be
reordered to have all its $\mathcal A$ factors first, so the image
is exactly the generated subcategory. The symmetry gives the
monoidal structure as in the universal property above.
\end{proof}

For example, the subcategories $\SVec_k$ and
$\operatorname{Ver}_p^+$ inside $\operatorname{Ver}_p$, for $p\geq5$,
have only the simple unit in common and generate the whole category.
The proposition recovers their product decomposition.

\subsection{Equivariant objects}

An action of a group $G$ on a category requires coherent functors,
not just a homomorphism into the group of isomorphism classes of
autoequivalences. Concretely, choose autoequivalences $T_g$ and
natural isomorphisms
\[
 \mu_{g,h}:T_gT_h\xrightarrow{\sim}T_{gh},
\]
with a unit identification for $T_1$, satisfying the associativity
cocycle identity. For a monoidal action, the $T_g$ and the
isomorphisms are monoidal.

\begin{definition}
An \emph{equivariant object} is an object $X$ together with
isomorphisms $u_g:T_g(X)\to X$ satisfying
\[
 u_{gh}\circ\mu_{g,h,X}=u_g\circ T_g(u_h)
\]
and the unit condition. A morphism $f:(X,u)\to(Y,v)$ satisfies
$f\circ u_g=v_g\circ T_g(f)$ for all $g$. The resulting category
is denoted $\mathcal C^G$.
\end{definition}

For a monoidal action, tensor products of the maps $u_g$ give
$\mathcal C^G$ a monoidal structure. If the action respects a
symmetry, the symmetry is equivariant as well. In the rest of this
lecture the action is \emph{trivial}: $T_g=\id$ and all coherence
maps are identities. Then an equivariant object is simply
$(X,\rho)$ with $\rho:G\to\Aut_{\mathcal C}(X)$ a group
homomorphism. For example, $\Vect_k^G=\Rep_k(G)$.

\begin{proposition}\label{prop:lec18-equivariant-product}
Let $\mathcal C$ be a semisimple symmetric tensor category and $G$
a finite group acting trivially. There is a symmetric equivalence
\[
 \mathcal C^G\simeq\mathcal C\boxtimes\Rep_k(G).
\]
\end{proposition}
\begin{proof}
Write an object of $\mathcal C$ as
$X=\bigoplus_S S\otimes_k M_S$, where the $S$ are simple and the
$M_S=\Hom(S,X)$ are finite-dimensional multiplicity spaces. The
notation $S\otimes_k M_S$ means a direct sum of copies of $S$.
Since $\End(S)=k$ and morphisms between distinct simples vanish, a
$G$-action on $X$ is exactly a $G$-representation on each $M_S$.
Equivariant morphisms are exactly the equivariant linear maps on
these spaces. This gives the claimed equivalence of categories.
It is the functor sending $X\boxtimes V$ to $X\otimes_k V$ with
$G$ acting on $V$, and the componentwise tensor actions show that
it is symmetric monoidal.
\end{proof}

\subsection{Semisimplifying equivariant objects}

\begin{lemma}
A simple object $S$ of a semisimple rigid symmetric category has
nonzero categorical dimension.
\end{lemma}
\begin{proof}
Rigidity gives
$\dim_k\Hom(\1,S\otimes S^*)=\dim_k\End(S)=1$, so $\1$ occurs
with multiplicity one in $S\otimes S^*$. Both coevaluation and
evaluation composed with the symmetry are nonzero, by the duality
identities. In a semisimple category, a nonzero inclusion and a
nonzero projection of a simple summand of multiplicity one have
nonzero composite. This composite is $\dim(S)$.
\end{proof}

This assertion concerns simple objects. It does not say that all
nonzero objects have nonzero dimension: in characteristic $p$,
$\1^{\oplus p}$ has dimension zero.

\begin{proposition}\label{prop:lec18-equivariant-semisimplification}
In the situation of Proposition~\ref{prop:lec18-equivariant-product},
\[
 \overline{\mathcal C^G}
 \simeq\mathcal C\boxtimes\overline{\Rep_k(G)}
\]
as symmetric tensor categories.
\end{proposition}
\begin{proof}
The indecomposable objects of $\mathcal C^G$ are $S\boxtimes V$,
where $S$ is simple in $\mathcal C$ and $V$ is an
\emph{indecomposable} representation of $G$; $V$ need not be simple.
The trace of an endomorphism of such an object is $\dim(S)$ times
its ordinary trace on $V$. Thus nilpotent endomorphisms have trace
zero, and
\[
 \dim(S\boxtimes V)=\dim(S)\dim_k(V).
\]
By the lemma, this vanishes precisely when $\dim_k(V)$ vanishes
in $k$. More generally, on each $S$-block the trace pairing on
morphism spaces is the ordinary trace pairing on $G$-modules
multiplied by the nonzero scalar $\dim(S)$. Its radical is therefore
unchanged. Quotienting by these radicals gives exactly
$\mathcal C\boxtimes\overline{\Rep_k(G)}$. Tensor products and
symmetries descend from the equivalence before quotienting.
\end{proof}

\subsection{The cyclic action on a tensor power}

Assume now that $\operatorname{char}k=p>0$ and that $\mathcal C$ is
semisimple. The symmetry supplies an $S_p$-action on $X^{\otimes p}$.
Restricting to the cycle $\sigma=(1\,2\,\cdots\,p)$ defines a
functor
\[
 P:\mathcal C\longrightarrow\mathcal C^{C_p},\qquad
 X\longmapsto(X^{\otimes p},\sigma),\qquad
 f\longmapsto f^{\otimes p}.
\]
Canonical rearrangement of tensor factors makes $P$ symmetric
monoidal. It is not additive. Compose it with semisimplification:
\[
 \mathrm{Fr}_0:\mathcal C\xrightarrow{P}\mathcal C^{C_p}
 \longrightarrow\overline{\mathcal C^{C_p}}
 \simeq\mathcal C\boxtimes\operatorname{Ver}_p.
\]

The essential point is that the mixed terms in a tensor power become
negligible. In the expansion of $(X\oplus Y)^{\otimes p}$, the
summands are indexed by words of length $p$ in the letters $X,Y$.
There are two constant words, giving $X^{\otimes p}$ and
$Y^{\otimes p}$. Every other word has a free orbit under cyclic
rotation because $p$ is prime. Choose a summand $W$ in such an
orbit. Using the cyclic action to identify its translates, the sum
of the orbit is the equivariant object $W\otimes_k k[C_p]$.
Under the product description above, it is $W\boxtimes J_p$ and
vanishes after semisimplification. Consequently the canonical
inclusions and projections induce
\[
 \mathrm{Fr}_0(X\oplus Y)
 \cong\mathrm{Fr}_0(X)\oplus\mathrm{Fr}_0(Y).
\]
The same orbit argument applied to morphisms proves additivity; we
give that proof in the next lecture. Scalars satisfy
\[
 \mathrm{Fr}_0(\lambda f)=\lambda^p\mathrm{Fr}_0(f).
\]

To express this as a $k$-linear functor, let $\mathcal C^{(1)}$ have
the same underlying additive symmetric category as $\mathcal C$,
but scalar action $\lambda\mathbin{\cdot_{(1)}}f=\lambda^p f$.
The field $k$ is perfect, so this is again a category over $k$.
Twist the whole target of $\mathrm{Fr}_0$ in this manner, and use
the fact that $\operatorname{Ver}_p$ is defined over $\mathbb F_p$
to identify its twist with itself. The resulting functor is
\[
 \mathrm{Fr}:\mathcal C\longrightarrow
 \mathcal C^{(1)}\boxtimes\operatorname{Ver}_p.
\]
This is the Frobenius functor of
\cite[Section 3.3]{OstrikVerlinde}. After additivity is established,
it is exact because every exact sequence in the source splits.

\begin{remark}
Semisimplicity is essential to the product description used in this
construction. For a general symmetric tensor category one cannot
obtain the correct Frobenius functor simply by taking the ordinary
semisimplification of $\mathcal C^{C_p}$. The later lectures use a
relative construction to treat that case. The questions of when a
category admits a fiber functor to $\operatorname{Ver}_p$, and why
$\operatorname{Ver}_p$ cannot be further compressed by a tensor
functor, will be addressed after studying this Frobenius functor.
\end{remark}
